\cotefile{29}
\scanpage{23}
% archive sheet number: 22
base $T$ ''. Lorsque $T$ est régulier, les notions 1\textsuperscript{\uline{o}} et 2\textsuperscript{\uline{o}} coïncident. Le théorème de spécialisation pour le groupe fondamental tame peut se démontrer, par la technique de SGA 1, quand on utilise la notion de revêtement modérément ramifié au sens large, relativement aux\unsure{diviseurs|diviseur} $D_{\bar{t}}$ \inserted{dans les $S_{\bar{t}}$} ($D$ étant à croisement normaux relativement à $T$, $S$ propre et lisse sur $T$).

\uline{Remarques} Evidemment, pour tout $S'$ objet de $\underline{\mathrm{Et}}_S$, on définit de façon évidente $\operatorname{Rev}_{\deleted{\underline{R}}}^{\inserted{\underline{R}}}(S')$, en considérant la restriction de $\underline{R}$ à $\underline{\mathrm{Et}}_{S'}\quad(\underline{\mathrm{Et}}_S)_{/S'}$. Pour $S'$ variable, on obtient ainsi une catégorie fibrée $\underline{\mathrm{Rev}}^{\underline{R}}$ sur $\underline{\mathrm{Et}}_S$, dont les fibres $\underline{\mathrm{Rev}}_{S'}^{\underline{R}}$ sont les $\operatorname{Rev}^{\underline{R}}(S')$, et un foncteur cartésien
\[
\underline{r}:\quad\underline{\mathrm{Rev}}^{\underline{R}}\qquad\underline{\mathrm{Rev}},
\tag*{\unsure{(*)|(\deleted{=}*)|(\deleted{*}=)}}
\]
où on désigne pour abréger par $\underline{\mathrm{Rev}}$ la catégorie \inserted{fibrée sur $\underline{\mathrm{Et}}_S$} des revêtements des différents objets $S'$ de $\underline{\mathrm{Et}}_S$. (NB on suppose que pour tout $S'$, le foncteur réalisation géométrique sur $\underline{C}_{S'}$ est à valeurs dans les revêtements de $S'$, et que le foncteur réalisation géométrique sur $\operatorname{Rev}^{\underline{R}}(S')$ est bien défini). Il faut insister sur le fait que pour nous, la ``donnée de ramification'' $\underline{R}$ est un objet de nature essentiellement auxiliaire, destiné à définir la catégorie fibrée $\underline{\mathrm{Rev}}^{\underline{R}}$ (qui est même un \uline{\unsure{c\deleted{h}a\deleted{a}hmp|c\deleted{a}h\deleted{h}amp}}) et le foncteur cartésien $\underline{r}$ de (*); c'est \unsure{(*)|(\deleted{=}*)|(\deleted{*}=)} qui est l'objet véritable de l'étude. On regardera donc deux données de ramification $\underline{R}$, $\underline{R}'$ comme essentiellement équivalentes, si elles définissent des foncteurs (*) équivalents (par une équivalence, bien sûr, induisant l'identité sur $\underline{\mathrm{Rev}}$). On peut se convaincre sur des exemples que $\underline{R}$ et $\underline{R}'$ peuvent avoir des allures très différentes. Ainsi, sous les conditions de l'exemple 3), si $\underline{R}'$ est défini par une sous-catégorie fibrée \deleted{de} $\underline{C}'$ de $\underline{C}$
